\section{A checked improvement beyond the incoming proposals: balanced Rocq certificates}

The incoming source-proof packets already contain the main semantic reductions:
ordinary rather than offset bisimulation for K215's zero-offset certificate,
a symmetry lemma eliminating a separately computed reverse check, and compact
positive-integer path codes. One Lean packet additionally identifies paths
with the same implicit zero extension, reducing the certificate from 2,240 to
1,426 checkpoints. These are prior results, not new contributions of this
experiment. The improvement implemented here is a verified balanced lookup
representation for the Rocq checker, which the incoming Rocq packet leaves as
future work.

\paragraph{Why the representation matters.}
The ordinary checker stores a list of related checkpoint pairs
\((P_i,Q_i,j_i,k_i)\). At a decrement checkpoint, it checks the two successor
segments against entries \(j_i\) and \(k_i\). The old lookup is
\texttt{nth\_error C (N.to\_nat j)}. Its traversal can inspect linearly many
list constructors on every lookup; over a certificate of \(n\) entries the
indexing component has a quadratic worst case. Neither the semantic
bisimulation theorem nor the segment interpreter requires this representation.

\paragraph{New definitions.}
A certificate tree is empty, a leaf containing one entry, or a node
\((c,L,R)\) carrying a binary cut \(c\). Its flattening is the in-order list
of entries. Validity asserts recursively that \(c\) equals the number of
entries in \(L\). Lookup at binary index \(j\) descends into \(L\) when
\(j<c\) and otherwise into \(R\) at \(j-c\). A leaf returns its entry
exactly at index zero. In particular, out-of-range indices are handled
correctly; no implicit totalization of list lookup is introduced.

\paragraph{Refinement theorem.}
For every valid certificate tree \(T\), every index \(j\), and every pair of
program decoder trees \(t_1,t_2\), the compiled development proves
\[
 \operatorname{lookup}(T,j)=
 \operatorname{nth\_error}(\operatorname{flatten}(T),\operatorname{toNat}(j))
\]
and
\[
 \operatorname{checkTree}(t_1,t_2,T)=
 \operatorname{check}(t_1,t_2,\operatorname{flatten}(T)).
\]
The first proof is structural induction on \(T\), using the two standard
list-append lookup cases and the cut invariant. The second proof transports
the lookup identity through the unchanged branch checks, entry checks, and
universal Boolean fold. It reuses the incoming one-sided ordinary
bisimulation theorem. No new semantic simulation argument, axiom, or admitted
goal is required.

\paragraph{Exact K215 instance.}
The concrete tree has 2,240 leaves and height 13 when the leaf is counted.
The kernel checks both its size invariant and
\[
 \operatorname{flatten}(T_{215})=\texttt{Plain215}.
\]
It then checks the complete forward certificate and proves the unchanged
all-register statement
\[
 \forall R,\quad
 \operatorname{rm\_halts}(p_{215},([],R))
 \ \longleftrightarrow\
 \operatorname{rm\_halts}(\mathrm{dagPerm},([],R)).
\]
The exact current \texttt{Prog215.v}, \texttt{DagPerm.v}, and original
certificate source match the incoming inputs byte for byte at repository
commit \texttt{53269bcab1efc0f82d3022645bd7cf0c22d84959}.
The new generic module is 194 lines and 6,480 bytes; its concrete data wrapper
and source theorem are 23 and 21 lines. This adds reusable infrastructure
and is not claimed to decrease the total source-line count.

\paragraph{Quantitative result and its limits.}
For the exact 2,239 decrement checkpoints and one halt checkpoint, there are
4,479 successor lookups including the initial lookup. List lookup traverses
5,018,706 constructors; balanced lookup traverses 54,514 tree nodes, a factor
of approximately 92.06. The independent census is reproducible. The formal
lemma \texttt{tree215\_lookup\_bound} proves the bound of 13 visited tree
nodes for every binary index, including invalid indices. This count concerns
only indexing; it is not a claim about total kernel work or elapsed time.
For balanced \(n\)-entry trees and \(O(n)\) queries, node visits are
\(O(n\log n)\). Charging comparisons and subtractions on binary indices
bit by bit gives \(O(n\log^2 n)\) indexing work. Program decoding and
register-segment computations remain unchanged.

One matched run in the shared environment recorded
\texttt{Time vm\_compute} wall times of 16.008 seconds for the plain checker
and 11.314 seconds for the balanced checker. Complete source-module compiler
times were 33.103 and 24.772 seconds. These are illustrative single-run
measurements. Moreover, the new tree-construction and exactness module took
13.159 seconds, so the experiment does not establish a first-build or full
endpoint speedup. A production generator could emit balanced tree data
directly and retain flattening equality as a separate audit target; that
extension has not been measured here.

\paragraph{Checked scope and integration.}
The new modules compile under Coq 8.18.0 with closed assumption reports.
The generic checker and its dependency closure additionally pass recursive
\texttt{coqchk}; the selected recursive audit does not include the concrete
source certificate. Dependencies were freshly compiled during this task,
then reused only within the task while iterating on the new module.
No full 237-source ZFC endpoint rebuild or Rocq 9.1 replay is claimed.
Production adoption consists of preserving the public source theorem name,
using the new proof body, updating the generator, and rebuilding
\texttt{Correct215} and \texttt{CorrectDense215} with fresh assumption reports.
The finite machine, source program, calculus, and endpoint statements are
unchanged. This work supplies no encoded PA or ZFC-internal derivation.
